\section{Policy Training} \label{sec:policy_training}

%TODO_PM_21_05: Possibly use the short version below and move the current text to the appendix.
% To train the policy (TRAIN\_RL in Algorithm \ref{alg:basic_loop}) we use the proximal policy optimization (PPO) algorithm~\cite{ppo}. It generates rollouts using the simulated environment $env'$ and uses them to improve the policy $\pi$. The fundamental difficulty lays in imperfections of the model compounding over time. To mitigate this problem we use short rollouts of $env'$ starting from ground-truth frames sampled from the buffer $D$. End of such rollout is not treated as end of episode in policy learning - the value estimate provided by PPO is used to mimic the following rewards. 

% We stress that imperfections in the model are unavoidable, as it is trained using only data generated with the current policy. For this reason in Algorithm~\ref{alg:basic_loop} we collect data iteratively, using the newest policy in each loop's COLLECT.

We will now describe the details of SimPLe, outlined in Algorithm \ref{alg:basic_loop}.  In step 6 we use the proximal policy optimization (PPO) algorithm~\citep{ppo} with $\gamma=0.95$. The algorithm generates rollouts in the simulated environment $env'$ and uses them to improve policy $\pi$. The fundamental difficulty lays in imperfections of the model compounding over time. To mitigate this problem we use short rollouts of $env'$. Typically every $N=50$ steps we uniformly sample the starting state from the ground-truth buffer $D$ and restart $env'$ (for experiments with the value
of $\gamma$ and $N$ see Section~\ref{subabl}). Using short rollouts may have a degrading effect as the PPO algorithm does not have a way to infer effects longer than the rollout length. To ease this problem, in the last step of a rollout we add to the reward the evaluation of the value function.
% Isn't this present in the standard PPO as well?
Training with multiple iterations re-starting from trajectories gathered in the real environment is  new to our knowledge. It was inspired by the classical Dyna-Q algorithm and, notably, in the Atari domain no comparable results have been achieved.

%We stress that imperfections in the model are unavoidable, as the model is trained using only data generated with the current policy. Moreover, the training of the world model is prone to overfitting as we try to gather as little data as possible from the real environment. Finally, sometimes the model fails in a catastrophic way by losing semantic coherence (e.g., a disappearing ball in Pong). 

The main loop in Algorithm \ref{dpll} is iterated $15$ times (cf. Section~\ref{subabl}). The world model is trained for $45$K steps in the first iteration and for $15$K steps in each of the following ones.  Shorter training in later iterations does not degrade the performance because the world model after first iteration captures already part of the game dynamics and only needs to be extended to novel situations.

In each of the iterations, the agent is trained inside the latest world model using PPO. In every PPO epoch we used 16 parallel agents collecting 25, 50 or 100 steps from the simulated environment $env'$ (see Section \ref{subabl} for ablations). The number of PPO epochs is $z\cdot 1000$, where $z$ equals to 1 in all passes except last one (where $z = 3$) and two passes number 8 and 12 (where $z = 2$). This gives $800$K$\cdot z$ interactions with the simulated environment in each of the loop passes. In the process of training the agent performs  $15.2$M interactions with the simulated environment $env'$.
